% sections/05_results.tex

\section{Results}
\label{sec:results}

The final seed-0 training runs and locked confirmation evaluations were still in
progress when this manuscript revision was prepared. We therefore do not promote the
five-case monitoring curves to quantitative results. Those curves are intended for
training diagnostics only and are especially unsuitable for comparison with earlier
runs, which used a different mask-loss normalization and a different five-case
monitoring cohort.

To make the development history auditable while the final evaluations are pending,
Fig.~\ref{fig:training-dynamics} records four corrected-run histories, including the
exploratory FiLM branch. The figure uses raw validation events without dashboard
smoothing and is a training diagnostic rather than evidence for the final method
ranking.

\begin{figure}[t]
  \centering
  \includegraphics[width=\textwidth]{training_dynamics}
  \caption{Interim EMA diagnostics for four seed-0 development runs on the same
  deterministic five-case monitoring cohort. The snapshot contains concat fixed,
  concat random, concat weighted, and exploratory FiLM weighted runs. Fixed-mask
  monitoring occurs every 5,000 steps and augmented-run monitoring every 10,000
  steps. Curves end at the latest available event and are not used for checkpoint or
  method selection. Final claims use the predeclared development and confirmation
  cohorts.}
  \label{fig:training-dynamics}
\end{figure}

The incomplete intermediate-checkpoint screen is not used for model selection or
reported as a final result. Ensemble selection uses saved per-case metrics on the
fixed 25-case development subset without dashboard smoothing.

\subsection{Development-Set Ensemble Selection}
\label{sec:results:ensemble-selection}
% Canonical per-case sources on Hendrix:
% results/paper300k-concat-fixed-s0-maxnorm-dev25-ensemble-screen-n5/ensemble_metrics/
% results/paper300k-concat-random-s0-maxnorm-dev25-ensemble-screen-n5/ensemble_metrics/
% results/paper300k-concat-weighted-s0-maxnorm-dev25-ensemble-screen-n5/ensemble_metrics/

Table~\ref{tab:ensemble-selection} reports the preconfirmation ensemble screen on
the 25-case development subset. For the fixed-mask arm, ensembling had no material
effect on the aggregate metrics: over all five candidates, SSIM spanned only
$0.709715$--$0.709722$, PSNR $18.178827$--$18.179194$ dB, and MSE
$0.015027$--$0.015028$ at the reported precision. Although distinct seeds changed
the reconstruction inside the scored healthy region for all 25 cases, averaging did
not consistently reduce the fixed arm's error.

The random and weighted augmentation arms behaved differently. Mean aggregation
with $N=5$ achieved the lowest joint rank for both arms, whereas $N=1$ ranked last.
The fixed arm slightly preferred mean aggregation with $N=3$ by rank, but its
aggregate metric differences were negligible and mean $N=5$ was worse than $N=1$.
Given that five trajectories require approximately five times as many sampler calls,
we freeze $N=1$ for the fixed arm and mean $N=5$ for the random and weighted arms.
The random and weighted runs also retain their first-trajectory $N=1$ outputs, so the
confirmation results can distinguish an equal-sampling comparison from the selected
arm-specific inference policies.

\begin{table}[b]
  \centering
  \caption{Challenge-style mean rank on the 25-case development subset for the
  terminal EMA checkpoints. Ranks are averaged over SSIM, PSNR, and MSE; lower is
  better. The $N=1$ mean/median duplicate is shown once. Bold values mark the best
  configuration within each column.}
  \label{tab:ensemble-selection}
  \small
  \setlength{\tabcolsep}{5pt}
  \begin{tabular}{@{}lrrrr@{}}
    \toprule
    Configuration & Fixed & Random & Weighted & Pooled \\
    \midrule
    $N=1$                 & 2.9200 & 4.7067 & 4.4000 & 4.0089 \\
    Mean, $N=3$           & \textbf{2.7867} & 2.1333 & 2.6933 & 2.5378 \\
    Median, $N=3$         & 2.9733 & 3.5333 & 3.7733 & 3.4267 \\
    Mean, $N=5$           & 3.2400 & \textbf{1.6000} & \textbf{1.5867}
                          & \textbf{2.1422} \\
    Median, $N=5$         & 3.0800 & 3.0267 & 2.5467 & 2.8844 \\
    \bottomrule
  \end{tabular}
\end{table}

% Activate these templates after the corresponding manifests have been audited.
% The checkpoint and ensemble figures should replace the interim training-dynamics
% figure if space is limited.
%
% \begin{figure}[t]
%   \centering
%   \includegraphics[width=\textwidth]{checkpoint_selection}
%   \caption{Checkpoint selection on the fixed 25-case development subset. Panels
%   show the candidate means for SSIM, PSNR, and MSE and the paired joint rank
%   obtained by averaging per-case ranks across the three challenge metrics. The
%   star marks the lowest-rank candidate. No dashboard smoothing is applied.}
%   \label{fig:checkpoint-selection}
% \end{figure}
%
% \begin{figure}[t]
%   \centering
%   \includegraphics[width=\textwidth]{ensemble_selection}
%   \caption{Ensemble selection on the same 25 development cases. Mean and median
%   aggregation are compared for one, three, and five seeded reconstructions per
%   case using the checkpoint-selection rank rule.}
%   \label{fig:ensemble-selection}
% \end{figure}

The final quantitative table will be populated from the immutable aggregate JSON and
per-case CSV files produced by the shared evaluator. Its primary columns will report
mean $\pm$ sample standard deviation for SSIM, PSNR, and MSE on the locked 75-case
confirmation cohort. Autoencoder and Pix2Pix baselines will be scored on the same
case manifests with the same official evaluator. The official validation submission
will be reported separately because its ground truth is hidden.

The qualitative figure will select the cases nearest the 10th, 50th, and 90th
percentiles of confirmation-cohort SSIM under the frozen method. Each row will show
the voided input, one-sample reconstruction, selected-ensemble reconstruction,
ground truth, and absolute-error maps. Case selection, slice selection, intensity
windows, and metric labels are deterministic and saved with the figure manifest.

% \begin{figure}[t]
%   \centering
%   \includegraphics[width=\textwidth]{qualitative_reconstructions}
%   \caption{Representative reconstructions from the locked confirmation cohort.
%   Rows are selected deterministically at the 10th, 50th, and 90th percentiles of
%   frozen-method SSIM. Cyan outlines the complete inpainting hole and dashed yellow
%   outlines the healthy region used for scoring. All anatomical panels in a row use
%   one intensity window; both error columns use one pooled, 99th-percentile-clipped
%   normalized color scale.}
%   \label{fig:qualitative-reconstructions}
% \end{figure}

% TODO: Insert final figures and tables only after their manifests verify the
% predeclared case files, full-DDPM sampling, EMA weights, and frozen selection.
